\documentclass[12pt,a4paper]{article}

% Compile with: tectonic Lab9_Moisei_Daniel_FAF241_Report.tex
\usepackage{fontspec}
\setmainfont{Liberation Serif}
\setmonofont{Liberation Mono}
\usepackage[left=2cm,right=1cm,top=2cm,bottom=2cm]{geometry}
\usepackage{setspace}
\usepackage{graphicx}
\usepackage[export]{adjustbox}
\usepackage{xcolor}
\usepackage{listings}
\usepackage{titlesec}
\usepackage{array}
\usepackage{tabularx}
\usepackage{caption}
\usepackage{float}
\usepackage{chngcntr}
\usepackage[hidelinks]{hyperref}

\graphicspath{{figures/}}
\counterwithin{figure}{section}
\definecolor{codebackground}{HTML}{F2F2F2}
\setlength{\parindent}{1.25cm}
\setlength{\parskip}{0pt}
\onehalfspacing
\setlength{\intextsep}{8pt plus 2pt minus 2pt}
\captionsetup[figure]{font=small,skip=4pt,justification=centering}

\titleformat{\section}
  {\normalfont\bfseries\centering\fontsize{13}{16}\selectfont}
  {\thesection}{0.8em}{\MakeUppercase}
\titleformat{name=\section,numberless}
  {\normalfont\bfseries\centering\fontsize{13}{16}\selectfont}
  {}{0pt}{\MakeUppercase}
\titleformat{\subsection}
  {\normalfont\bfseries\fontsize{12}{15}\selectfont}
  {\thesubsection}{0.7em}{}
\titlespacing*{\section}{0pt}{0pt}{12pt}
\titlespacing*{\subsection}{0pt}{10pt}{6pt}

\lstset{
  basicstyle=\ttfamily\small,
  backgroundcolor=\color{codebackground},
  frame=none,
  breaklines=true,
  columns=fullflexible,
  keepspaces=true,
  showstringspaces=false,
  keywordstyle=\ttfamily,
  commentstyle=\ttfamily,
  stringstyle=\ttfamily,
  xleftmargin=0.4cm,
  xrightmargin=0.2cm,
  aboveskip=6pt,
  belowskip=10pt
}
\newcommand{\taskfigure}[4][]{%
  \begin{figure}[H]
    \centering
    \includegraphics[scale=0.52,max width=\textwidth,max height=0.75\textheight]{figures/#2}
    \caption{#3}
    \label{fig:#4}
  \end{figure}
}
\newcommand{\inlinecmd}[1]{\texttt{\detokenize{#1}}}

\begin{document}

\begin{titlepage}
  \begin{center}
    \includegraphics[width=3.4cm]{figures/LOGO_UTM.jpg}\par
    \vspace{0.25cm}

    {\bfseries\fontsize{14}{17}\selectfont
      MINISTRY OF EDUCATION, CULTURE AND RESEARCH\par
      OF THE REPUBLIC OF MOLDOVA\par
      Technical University of Moldova\par
      Faculty of Computers, Informatics and Microelectronics\par
      Department of Software Engineering and Automation\par
    }

    \vspace{1.25cm}
    {\bfseries\fontsize{16}{19}\selectfont Daniel Moisei FAF-241\par}

    \vspace{1.35cm}
    {\bfseries\fontsize{34}{39}\selectfont Report\par}
    \vspace{0.2cm}
    {\itshape\fontsize{20}{24}\selectfont Laboratory Work No. 9\par}
    {\bfseries\itshape\fontsize{20}{24}\selectfont Operating Systems\par}

    \vspace{0.8cm}
    {\fontsize{18}{22}\selectfont Topic: Watching System Calls with strace\par}
  \end{center}

  \vfill
  \begin{flushright}
    {\fontsize{15}{18}\selectfont Checked by:}\par
    \textbf{Reițman Patricia}, \textit{university assistant}
  \end{flushright}

  \vfill
  \begin{center}
    {\bfseries\fontsize{14}{17}\selectfont Chișinău - 2026}
  \end{center}
\end{titlepage}

\pagenumbering{arabic}
\setcounter{page}{2}
\tableofcontents
\newpage

\section*{Introduction}
\addcontentsline{toc}{section}{Introduction}

The purpose of this laboratory work was to observe the boundary between a program and the Linux kernel: the system call. The tool \texttt{strace} runs a program and prints every system call it makes, with its arguments and its return value. I used it first on a familiar command, \texttt{ls -l}, then on a small C program, \texttt{syscalls.c}, that opens, writes, reads and closes a file directly through system calls, and finally to count how many calls of each kind a program makes.

The work was done on CachyOS (Arch Linux). \texttt{strace} was not installed by default; it was installed with the AUR helper \inlinecmd{paru -S strace} instead of \texttt{apt}, which the handout mentions for Debian and Ubuntu.

\newpage
\section{See the Calls of a Real Command}

\subsection{Tracing ls -l}

The command \inlinecmd{strace -f -e trace=openat,read,write,close ls -l} prints only the four file-related calls. The beginning of the trace is a long burst of \texttt{openat}, \texttt{read} and \texttt{close}: the dynamic loader opens \texttt{/etc/ld.so.cache} and the shared libraries (\texttt{libcap.so.2}, \texttt{libc.so.6}), then \texttt{ls} opens the locale files, the current directory, and \texttt{/etc/passwd} and \texttt{/etc/group} to translate the numeric owner and group into names. Several \texttt{openat} calls return \texttt{-1 ENOENT}: these are translation files that do not exist for the English locale, and \texttt{ls} simply continues.

Figure~\ref{fig:ls-trace} shows the end of the trace. After all the files have been read and closed, the whole listing is printed by a single call, \texttt{write(1, "total 4\textbackslash n...", 66) = 66}, and the program ends by closing descriptors 1 and 2.

\taskfigure{part1-ls-trace.png}{The end of the trace of ls -l: one write to descriptor 1}{ls-trace}

\subsection{How many system calls}

To count all the system calls, not only the four traced above, the full trace was filtered with \texttt{grep}, counting the lines that start with a call name (Figure~\ref{fig:ls-count}). One simple \texttt{ls -l} made 259 system calls. Most of them are start-up work (\texttt{execve}, \texttt{mmap}, \texttt{mprotect}, \texttt{brk}, loading of the libraries) and user-name lookups; the real output is a single \texttt{write}.

\taskfigure{part1-ls-count.png}{The number of system calls made by ls -l}{ls-count}

\subsection{Answer}

\texttt{ls} prints to the screen with the system call \texttt{write}, and it writes to descriptor \textbf{1}, the standard output, which is connected to the terminal. The listing of the directory is sent in one \texttt{write} call of 66 bytes (the directory contained only \texttt{syscalls.c}).

\newpage
\section{A Program That Calls the Kernel Directly}

\subsection{The program}

The program \texttt{syscalls.c} opens \texttt{note.txt}, writes one line into it, closes it, then opens it again and reads the line back. The file itself is accessed only through the system-call wrappers \texttt{open}, \texttt{write}, \texttt{read} and \texttt{close}; \texttt{printf} is used only to report what happened.

\lstinputlisting[language=C,caption={Implementation of syscalls.c}]{../syscalls.c}

The program compiled without warnings and printed the expected output (Figure~\ref{fig:syscalls-run}): descriptor 3, 14 bytes written, and the same 14 bytes read back.

\taskfigure{part2-syscalls-run.png}{Compilation and execution of syscalls}{syscalls-run}

\subsection{The program seen from outside}

Under \inlinecmd{strace -e trace=openat,write,read,close ./syscalls} the C calls of the program appear one-to-one as system calls (Figure~\ref{fig:syscalls-strace}). Before them there are a few set-up calls of the dynamic loader, which opens \texttt{ld.so.cache} and \texttt{libc.so.6}. Then come the six calls of the program:

\begin{lstlisting}
openat(AT_FDCWD, "note.txt", O_WRONLY|O_CREAT|O_TRUNC, 0644) = 3
write(3, "hello, kernel\n", 14)         = 14
close(3)                                = 0
openat(AT_FDCWD, "note.txt", O_RDONLY)  = 3
read(3, "hello, kernel\n", 63)          = 14
close(3)                                = 0
\end{lstlisting}

The C function \texttt{open} is implemented in glibc with the system call \texttt{openat}, using \texttt{AT\_FDCWD} to mean ``relative to the current directory''. The \texttt{read} asks for 63 bytes (\texttt{sizeof(buf) - 1}) but receives 14, because the file contains only 14 bytes.

Between these calls the trace also shows the three \texttt{printf} lines, each as its own \texttt{write} to descriptor 1:

\begin{lstlisting}
write(1, "open returned descriptor 3\n", 27) = 27
write(1, "write sent 14 bytes\n", 20)        = 20
write(1, "read got 14 bytes: hello, kernel"..., 33) = 33
\end{lstlisting}

Each one comes right after the \texttt{printf} that produced it, because the standard output is a terminal and the C library sends it line by line. The text of the program appears on the screen in the middle of the \texttt{strace} line, since both write to the same terminal: the program to descriptor 1 and \texttt{strace} to descriptor 2.

\taskfigure{part2-syscalls-strace.png}{The system calls of syscalls traced by strace}{syscalls-strace}

\subsection{Answer}

Descriptors 0, 1 and 2 are already taken when the program starts: every process inherits from the shell its \textbf{standard input} (0, the keyboard), its \textbf{standard output} (1, the screen) and its \textbf{standard error} (2, the screen). The kernel always returns the smallest free descriptor, so the first file opened by the program gets 3. After \texttt{close(3)} the number is free again, which is why the second \texttt{openat} also returns 3.

\newpage
\section{Count One Kind of Call}

\subsection{A summary of all the calls}

The option \texttt{-c} makes \texttt{strace} print, instead of every call, a table with the number of times each system call was made and how many of them failed (Figure~\ref{fig:syscalls-count}). The program's own output was sent to \texttt{/dev/null} so that only the table remains.

\taskfigure{part3-syscalls-count.png}{The summary of the system calls of syscalls}{syscalls-count}

\texttt{syscalls} made 41 system calls in total. Its own work is visible in the table: 4 \texttt{openat} (2 for the loader, 2 for \texttt{note.txt}), 2 \texttt{read} (the library header and the 14 bytes of the file), 2 \texttt{write} (the 14 bytes into \texttt{note.txt}, and a single write with all three \texttt{printf} lines: with the output sent to \texttt{/dev/null} instead of the terminal, the C library buffers the lines and writes them together at exit) and 4 \texttt{close}. The 2 errors are normal start-up probes: \texttt{access} checks for a preload file that does not exist, and \texttt{ioctl} checks whether the standard output is a terminal, which fails for \texttt{/dev/null}. All the other rows (\texttt{mmap}, \texttt{mprotect}, \texttt{brk}, \texttt{execve}, \ldots) belong to the dynamic loader and the C library, not to the code of the program.

\subsection{Counting only the write calls of ls}

To count one kind of call, the trace was restricted to \texttt{write} with \inlinecmd{-e trace=write} and the lines were counted (Figure~\ref{fig:ls-writes}). Out of its 259 system calls, \texttt{ls -l} makes exactly \textbf{one} \texttt{write}: the C library collects the whole listing in its buffer and hands it to the kernel at once. Each system call has a cost, so buffering the output saves many crossings between user mode and kernel mode.

\taskfigure{part3-ls-writes.png}{Only one write call in ls -l}{ls-writes}

\newpage
\section*{Conclusion}
\addcontentsline{toc}{section}{Conclusion}

This laboratory work made the system call visible. Every action that touches something outside the process, such as opening a file, reading a directory or printing to the screen, is a request to the kernel, and \texttt{strace} shows each request with its arguments and result. Even a small command like \texttt{ls -l} made 259 system calls, most of them to load the shared libraries and prepare the program, and only one to print the result.

The small program \texttt{syscalls.c} showed that the C functions \texttt{open}, \texttt{write}, \texttt{read} and \texttt{close} map directly onto system calls with the same arguments, with \texttt{open} implemented as \texttt{openat}. The descriptor numbers confirmed the rule seen in Laboratory Work 7: descriptors 0, 1 and 2 are inherited from the shell, and the kernel always gives the smallest free number. Finally, \texttt{strace -c} showed how the C library reduces the number of system calls by buffering output, so that several \texttt{printf} lines, or a whole directory listing, leave the program in a single \texttt{write}.

\end{document}
